\medskip 
{\color{gray}\hrule}
\section{Factoring Primes in a Monogenic Number Field}
\medskip 
{\color{gray}\hrule}
\medskip 

\begin{flushleft}
    \quad\textbf{Theorem (Dedekind's Theorem):} Let K = Q($\theta$) be an algebraic number field of degree n, such that 
    \[
    \mathcal{O}_K = \mathbb{Z} + \mathbb{Z}\theta + \cdots + \mathbb{Z}\theta^{n-1}
    \]
     (i.e. it's a monogenic number field). 
    Let $f(x) \in \mathbb{Z}[x]$ be the minimal polynomial of $\theta$, and let $p$ be a prime number such that
    \[
    p \nmid [\mathcal{O}_K : \mathbb{Z}[\theta]].
    \]

    \medskip
    
     Define the reduction modulo $p$ map
     \[
     \phi : \mathbb{Z}[x] \to \mathbb{Z}_p[x] .
     \]
     
     \medskip
     Then for $f(x) \in \mathbb{Z}[x]$, we can write
     \[
     \phi(f(x)) = g_1^{e_1} \cdots g_r^{e_r},
     \]
     where each $g_i(x) \in \mathbb{Z}_p[x]$ is distinct, monic, irreducible and each $e_i \in \mathbb{Z}$.
     
    Define $f_i(x)$ to be a monic polynomial of $\mathbb{Z}[x]$ such that $\phi(f_i(x)) = g_i(x)$.
    Set 
    \[
    Pi = <p, f_i(\theta)>, i = 1,2,...,r 
    \]
    to be the principal ideal generated by $p$ and $f_i(\theta)$. Each of these $P_i$'s are distinct
    prime ideals of $\mathcal{O}_K$ where the ideal generated by $p$ is equal to a product of these ideals. 
    In other words, we can write each of these ideals as follows
    \[
    <p>=P_1^{e_1} \cdots P_r^{e_r}.
    \]
    Define the norm of $P_i$ to be $p$ raised to the power of the degree of the corresponding $f_i$. We denote it as follows
    \[
    N(P_i)=p^{deg f_i}, i = 1,2,...,r.
    \]
     
    \textit{Proof}
     
    \medskip 

    The proof is broken down into 3 steps. First, we have to show that each $P_i = \langle p, f_i(\theta)\rangle$.
     
    \medskip 

    Let $\theta_i$ be a root of $g_i$, an irreducible polynomial of $\mathbb{Z}_p[x]$.
    Namely, the finite field 
    \[
    \mathbb{Z}_p[\theta_i] \cong \mathbb{Z}_p[x]/\langle g_i(x) \rangle. 
    \]
    (which holds by applying the first isomorphism theorem). 

    Let $v_i$: $\mathbb{Z}[\theta] \to \mathbb{Z}_i[\theta_i]$ denote a (surjective) homomorphism, sending $v_i(h(\theta)) = \bar{h}(\theta_i)$. 
    (the $\bar{h}(\theta_i)$ are from the map stated in the theorem's statement, it is a product of distinct, irreducible, monic polynomials in $\mathbb{Z}_p[x]$)
    In particular, we have that 
    \[ 
    \mathbb{Z}[\theta]/Ker(v_i) \cong v_i(\mathbb{Z}[\theta]) = \mathbb{Z}_p[\theta_i].
    \] 
    So, $Ker(v_i)$ is a prime ideal of $\mathbb{Z}[\theta] = \mathcal{O}_K$. We have that 
    \[
    v_i(p) = 0,\quad v_i(f_i(\theta)) > 0,
    \]
    and
    \[
    \bar{f_i}(\theta_i) = g_i(\theta_i) = 0,
    \]
    where these are polynomials in $\mathbb{F}_p[x]$.
    
    Therefore, $p$ and $\bar{f_i}(\theta) \in Ker(v_i)$. This implies that $\langle p, f_i(\theta) \rangle$ $\subseteq Ker(v_i)$ where $\langle p, f_i(\theta) \rangle$ is the principal ideal generated by $p$ and $f_i(\theta)$.

    For the reverse inclusion, let $g_i(\theta) \in Ker(v_i)$. Then, 
    \[
    \bar{g_i}(\theta_i) = v_i(g(\theta)) = 0, \text{ with } g_i \mid \bar{g}(x) \in \mathbb{F}_p[x].
    \]
 
    We have that 
    \[
    \bar{g}(x) = \bar{f_i}(x)\bar{h}(x), \text{ for some } \bar{h}(x) \in \mathbb{F}_p[x].
    \]
    
    

    By rearranging the equation, this results in 
    \[ 
    g-f_ih \in \mathbb{Z}[x]
    \] 
    with coefficients divisible by p (since this equation is 0 mod p).

    So, 
    \[ 
    g(\theta) = (g(\theta) - f_i(\theta)h(\theta)) + f_i(\theta)h(\theta)
    \]
    with the first term being in $\langle p \rangle$ and the second term being in $\langle f_i(\theta) \rangle$. 
    Therefore, $g(\theta)$ is a principal ideal generated by $\langle p, f_i(\theta) \rangle$, showing that $Ker(v_i) \subseteq \langle p, f_i(\theta) \rangle)$. 

    This proves that each ideal 
    \[ 
    P_i = \langle p, f_i{\theta} \rangle = Ker(v_i)
    \]
    are prime ideals of $\mathcal{O}_K$, for all $i = 1, ..., r$. 

    Next, we must prove that each $P_i$ are distinct. 
    

    Suppose that $P_i = P_j$, for some i and j $=1,...,r$. Then, 
    \[
    P_i = \langle p, f_i(\theta) \rangle = \langle p, f_j(\theta) \rangle = P_j.
    \]
    
    Let $f_j(\theta) \in P_j$. Since it is also generated by $p$ and $f_i$, then 
    \[
    f_j(\theta) = pg(\theta) + f_i(\theta)h(\theta), \text{ for } g, h \in \mathbb{Z}[x].
    \] 

    Now, we apply the map $v_i$ defined above on $f_j$. We get
    \[
    g_j(\theta_i) = \bar{f_j}(\theta_i) \bar{f_i}(\theta_i)\bar{h}(\theta_i) = g_i(\theta_i)\bar{h}(\theta_i) = 0.
    \]
    From this, we see that $g_i \mid g_j(x)$ in $\mathbb{Z}_p[x]$. So, $g_j{x}$ can be written as a product of $g_i(x)$ and $l(x)$, where $l(x)$ is a polynomial in $\mathbb{Z}_p[x]$. 
    However, since $g_i$ and $g_j$ are monic and irreducible polynomials in $\mathbb{Z}_p[x]$, the only choice for $l(x)$ is to be equal to 1. 

    This implies that $g_i = g_j$, and therefore $i=j$. Therefore, each prime ideal $P_i$ is distinct, for $i=1,...r$.

    Finally, we must prove that the principal ideal can be factored into a product of distinct prime ideals as follows
    \[
    \langle p \rangle = {P_1}^{e_1}...{P_r}^{e_r},
    \] 
    and that the norm of each $P_i$ is equal to $p^{deg(f_i)}$.

    To begin, we need the following lemma.
    
    \quad\textbf{Lemma:} Let $A$, $B_1$, and $B_2$ be ideals of a ring $R$. Then, we have 
    \[
    (A+B_1)(A+B2) \subseteq A+B_1B_2.
    \] We omit the proof of the lemma. 


    So, for ${P_1}^{e_1} \cdots {P_r}^{e_r}$ we have that:
    ${P_1}^{e_i} \cdots {P_r}^{e_r} = {\langle p, f_1(\theta) \rangle}^{e_1} \cdots {\langle p, f_r(\theta) \rangle}^{e_r}
    \subseteq \langle p \rangle + {\langle f_1(\theta) \rangle}^{e_1}...{\langle f_r(\theta) \rangle}^{e_r}$ (factoring out the p)
    $= \langle p \rangle + \langle f(\theta) \rangle = \langle p \rangle$. (since f is the minimal polynomial of theta - evaluates to 0)
    So, $\langle p \rangle \mid {P_1}^{k_1}...{P_r}^{k_r}$., for $k_i \in \{1,2, \dots, e_i\}, i=1,2,...r$. 
    Implying that $P_i = \langle p, f_i(\theta) \rangle \subseteq \langle p \rangle$. 

    We also have that $  {P_1}^{e_1}...{P_r}^{e_r} \mid\langle p \rangle$. Therefore, $\langle p \rangle = {P_1}^{e_1}...{P_r}^{e_r}$.  

    To prove that $N(P_i) = p^{(deg(f_i))}$, start with $\mathcal{O}_K/P_i = \mathbb{Z}[\theta]/P_i = \mathbb{Z}[\theta]/Ker(v_i) \cong v_i(\mathbb{Z}[\theta]) 
    = \mathbb{Z}_p[\theta_i]$. 

    Such that $N(P_i) = |\mathcal{O}_K/P_i| = |\mathbb{Z}_p[\theta_i]| = p^{d_i}$, where $d_i = deg(g_i) = deg(\bar{f_i})$.

    We then have that $p^n = N(\langle p \rangle) = N({P_1}^{k_1}...{P_r}^{k_r})$, by the multiplicative property of the norm, this is equal to
    $N(P_1)^{k_1} \dots N(P_r)^{k_r} = (p^{d_1})^{k_1} \dots (p^{d_r})^{k_r} = p^{d_1k_1+} \dots + p^{d_rk_r}$, with $d_1k_1 + \dots + d_rk_r = n$. 

    So, $\bar{f}(x) = \bar(f_1)(x)^{e_1} \dots \bar(f_r)(x)^{e_r}$, where $d_1e_1 + \dots + d_re_r = n$ since each $k_i \in \{1,2, \dots, e_i\}$.

    Therefore, we have that $k_i = e_i$. 
  
    So indeed, $\langle p \rangle = {P_1}^{e_1}...{P_r}^{e_r}$, and $N(P_i) = p^{d_i} = p^{deg(\bar{f_i})} = p^{deg(f_i)}$, for $i = 1, 2, \dots, r$. 

    Thus, the theorem is proved. $\square$

    
\end{flushleft}



{\color{gray}\hrule}
\subsection{Nom de la sous-section}

